\documentclass[aspectratio=169]{beamer}
\usetheme{metropolis}
\usepackage{graphicx}

\title{Is the wall behaving as expected?}
\subtitle{Grey-box monitoring of station 02 (2018--2026)}
\date{}

\begin{document}

\maketitle

\begin{frame}{The Problem}
    \textbf{Question:} Given the current environment, is the wall leaning exactly as expected, or doing something unusual?
    \vspace{0.5cm}
    \begin{itemize}
        \item Readings every 20 minutes from the medieval wall of Gubbio:
        \begin{itemize}
            \item Inclination (leaning)
            \item Air temperature
            \item Humidity
            \item Solar radiation
        \end{itemize}
        \item \textbf{Goal:} Model normal behavior and monitor for unexpected movements over 8 years.
    \end{itemize}
\end{frame}

\begin{frame}{The Dataset}
    \begin{itemize}
        \item We use three sets of environmental drivers to ensure reliability when the wall's own sensors (\texttt{str}) fail.
        \item External proxies: Local ground station (\texttt{gs}) and ERA5 reanalysis (\texttt{era5}).
    \end{itemize}
    \centering
    \includegraphics[width=\textwidth,height=0.6\textheight,keepaspectratio]{../outputs/GM_F01_regressor_sets.png}
\end{frame}

\begin{frame}{How We Model the Wall}
    \begin{itemize}
        \item We use a ``grey-box'' model (\textit{NeuralProphet}) separating inclination into:
        \begin{itemize}
            \item \textbf{Trend:} Long-term wandering.
            \item \textbf{Seasonality:} Annual and daily repeating cycles.
            \item \textbf{Drivers:} Temperature, humidity, radiation.
        \end{itemize}
    \end{itemize}
    \centering
    \includegraphics[width=\textwidth,height=0.5\textheight,keepaspectratio]{../outputs/GM_F06_decomposition.png}
\end{frame}

\begin{frame}{What Drives the Wall's Movement?}
    \begin{columns}
        \begin{column}{0.5\textwidth}
            \begin{itemize}
                \item \textbf{Slow Trend ($\approx$60\%):} Major wandering over years. Not a straight drift; it reverses.
                \item \textbf{Air Temp (15--24\%):} The wall and air heat together from solar radiation. The daily swing is captured instantaneously by air temperature.
                \item \textbf{Annual Cycle (12--17\%):} Delayed seasonal response.
            \end{itemize}
        \end{column}
        \begin{column}{0.5\textwidth}
            \centering
            \includegraphics[width=\textwidth,keepaspectratio]{../outputs/GM_F04_trend.png} \\
            \vspace{0.2cm}
            \includegraphics[width=\textwidth,keepaspectratio]{../outputs/GM_F05_seasonality.png}
        \end{column}
    \end{columns}
\end{frame}

\begin{frame}{Monitoring for Anomalies}
    \begin{itemize}
        \item \textbf{Rolling Expectation:} The model is re-trained every 30 days to establish a moving baseline.
    \end{itemize}
    \centering
    \includegraphics[width=0.8\textwidth,height=0.6\textheight,keepaspectratio]{../outputs/GM_F08_observed_expected.png}
\end{frame}

\begin{frame}{The Control Charts}
    We use three control charts to catch different anomalies:
    \begin{enumerate}
        \item \textbf{Fast Chart:} Catches sudden steps or spikes (e.g., instrument faults).
        \item \textbf{Daily Chart:} Catches changes in daily swing amplitude or phase.
        \item \textbf{Slow Chart:} Catches slow continuous drift over weeks.
    \end{enumerate}
    \vspace{0.3cm}
    \centering
    \includegraphics[width=0.9\textwidth,height=0.4\textheight,keepaspectratio]{../outputs/GM_F09_fast_chart.png}
\end{frame}

\begin{frame}{Behavior Across Outages}
    \begin{itemize}
        \item What happens when sensors go offline? The monitor flags when the wall resumes at an unexpected level.
    \end{itemize}
    \centering
    \includegraphics[width=0.8\textwidth,height=0.6\textheight,keepaspectratio]{../outputs/GM_F12_outage_bridges.png}
\end{frame}

\begin{frame}{Conclusion}
    \begin{itemize}
        \item The wall at station 02 acts largely as a thermometer, driven primarily by solar heating (captured via air temperature) and a slow, reversible structural drift.
        \item By separating this expected environmental response from raw data, our model and control charts provide a sensitive, automated way to monitor the true structural health.
    \end{itemize}
\end{frame}

\end{document}
